% Lösungen Einheit 4 von 5 – Zylinder (zusammenbau v0.1)
\einheitenkopf{Einheit 4 von 5 · Zylinder}

\erg{22}{a) cm²}
%% TODO Lösung der Erklärzeile 23g) fehlt in der Bank
\erg{23}{a) Durchmesser; Radius $4$ cm \quad b) $\pi \cdot 5^2 \cdot 8 \approx 628{,}3$ cm³ \quad c) $\pi \cdot 2^2 \cdot 9 \approx 113{,}1$ cm³ \quad d) $\pi \cdot 6^2 \cdot 4 \approx 452{,}4$ cm³ \quad e) $\pi \cdot 7^2 \cdot 5 \approx 769{,}7$ cm³ \quad f) $\pi \cdot 4^2 \cdot 6 \approx 301{,}6$ cm³ \quad g) \ldots \quad h) $r = 4$ cm; $\pi \cdot 4^2 \cdot 5 \approx 251{,}3$ cm³ \quad i) $\pi \cdot 2{,}5^2 \cdot 7{,}5 \approx 147{,}3$ cm³ \quad j) $\pi \cdot 11^2 \cdot 14 \approx 5\,321{,}9$ cm³ $\approx 5{,}3$ l \quad k) $2 \cdot \pi \cdot 4 \cdot 9 \approx 226{,}2$ cm² \quad l) $2 \cdot \pi \cdot 5^2 + 2 \cdot \pi \cdot 5 \cdot 9 \approx 439{,}8$ cm² \quad m) Umfang $2 \cdot \pi \cdot 2 \approx 12{,}6$ cm, also $5$ cm × $12{,}6$ cm \quad n) $500 : (\pi \cdot 5^2) \approx 6{,}4$ cm \quad o) $\sqrt{600 : (\pi \cdot 9)} \approx 4{,}6$ cm (erst teilen, dann die Wurzel) \quad p) $\pi \cdot 3{,}5^2 \cdot 9 \approx 346{,}4$ cm³ $\approx 346$ ml – das sind etwa $350$ ml \quad q) A: $\pi \cdot 3^2 \cdot 5 = \pi \cdot 45$; B: $\pi \cdot 6^2 \cdot 5 = \pi \cdot 180$; $180 : 45 = 4$ – vierfach, weil der Radius im Quadrat steht \quad r) $2{,}5$ l $= 2\,500$ cm³; Höhe $= 2\,500 : (\pi \cdot 7^2) \approx 16{,}2$ cm}
\erg{24}{a) Max hat den Durchmesser statt des Radius eingesetzt. Richtig: $r = 0{,}35$ m, $V = \pi \cdot 0{,}35^2 \cdot 0{,}9 \approx 0{,}35$ m³ \quad b) Der Radius steht im Quadrat: Der doppelte Radius gibt eine viermal so große Grundfläche, die Höhe bleibt gleich. \quad c) $\pi \cdot 28^2 \cdot 90 \approx 221\,671$ cm³ $\approx 221{,}7$ l; zwei Tage brauchen $280$ l – nein \quad d) Rechteck $4$ cm hoch und $2 \cdot \pi \cdot 1{,}5 \approx 9{,}4$ cm lang; oben und unten je ein Kreis mit $1{,}5$ cm Radius}

24d)

\netzzylinder{1.5}{4}{$1{,}5$ cm}{$4$ cm}

